\section{Package options}
\label{sec:package-options}

The package uses a key--value interface based on \pkg{l3keys2e}. %
Options may be passed at package load time or via \refCom{liquidsetup} in
the preamble or the document. %

\begin{blockwarning}
  See the runtime chapter for additions in version 0.3.0. Some keys are restricted to preamble usage only. %
  When this is the case, the description of the key is marked with a
  \preambleonly{} symbol. %
\end{blockwarning}

\subsection{General options}
\label{ssec:general-options}

These options are used for the general configuration of \pkg{liquidtex}. %

\begin{docKey}{verbose}{=\docValue{true}\textbar\docValue{false}}{default:
    \docValue{true}, initial: \docValue{false}}%
  \index{verbose!configuration}
  \index{logging!configuration}
  Enables verbose diagnostic messages.
\end{docKey}
\begin{docKey}{no-cache}{=\docValue{true}\textbar\docValue{false}}{default:
    \docValue{true}, initial: \docValue{false}}%
  \index{cache!deactivation}
  Disables the use of cache-rendered output.%
  Forces every request to render, bypassing persistent cached output.
\end{docKey}
\begin{docKey}[][doc description=\preambleonly]{outdir}{=\meta{path}}{default: \meta{path}}
  \index{output directory!configuration}
   Directory used to store rendered templates and cache files. %
\end{docKey}
\begin{docKey}[][doc description=\preambleonly]{use-auxdir}{=\docValue{true}\textbar\docValue{false}}{default:
    \docValue{true}, initial: \docValue{true}}%
  Indicates whether \refKey{outdir} should be interpreted relative to the
  engine’s auxiliary directory (e.g.\ under \meta{--auxdir} or
  \meta{--output-directory}). %
  This applies to relative cache and debug directories.
\end{docKey}

%\PrintChanges

\subsection{Template data}
\label{ssec:template-data-options}

\begin{docKey}{data}{=\meta{path}}{} %
  Sets the default data source used when rendering templates. %
  No default value is set, static templates may be rendered with no selected data source. %
\end{docKey}
\begin{docKey}[][doc new=2025-11-21]{data-sources}{=\meta{options}}{} %
  Adds or replaces individual data source aliases without clearing other sources. %
  The \meta{options} are passed as a key-value collection where the key is
  the alias of the data source and the key is the path to the file. %

  The data sourced passed with \refKey{data} is stored as
  \docValue{default} and can be overriden. %
\end{docKey}

\subsection{Engine selection}
\label{ssec:engine-selection-options}

\pkg{liquidtex} can render
Liquid templates with different backends. %
The rendering engine can be selected with: %

\begin{docKey}[][doc description=\preambleonly]{engine}{=\docValue{lua}\textbar\docValue{external}}{initial:
    \docValue{external}}
  Chooses the rendering backend.

  \begin{itemize}
  \item the \docValue{external} choice allows the user to choose an external
    engine and render from outside;
  \item the \docValue{lua} choice uses an internal engine to render the
    template.
  \end{itemize}

  More information on engine configuration is available in
  \cref{sec:engine-configuration}. %
\end{docKey}

%%% Local Variables:
%%% mode: LaTeX
%%% TeX-master: "../liquidtex"
%%% End:
